%=======================================================================
%  preambulo.tex — paquetes, estilos y macros
%
%  Recortado para el miniproyecto: sin pgfplots (las gráficas entran como
%  imágenes desde figuras/), sin algoritmos y sin entornos tipo teorema.
%=======================================================================

%----------------------- Codificación e idioma -------------------------
\usepackage[T1]{fontenc}
\usepackage[utf8]{inputenc}

% El ORDEN de las opciones importa y es la causa número uno de que un documento
% en español no compile: las opciones es-* las declara spanish.ldf, así que
% 'spanish' tiene que ir ANTES que ellas.  Con main=spanish babel pospone la
% carga del español y aborta con «Unknown option 'es-tabla'»; y como pdflatex
% termina con error, latexmk nunca llega a ejecutar biber y el PDF sale con
% todas las citas en [?].  El idioma principal es el último declarado.
%   es-tabla          -> «Tabla» en vez de «Cuadro»
%   es-nodecimaldot   -> no toca el punto decimal en modo matemático
%   es-noshorthands   -> desactiva los atajos < > " que rompen las flechas de TikZ
\usepackage[english,spanish,es-tabla,es-nodecimaldot,es-noshorthands]{babel}
\usepackage{csquotes}          % biblatex lo exige cuando hay babel

\usepackage{lmodern}
\usepackage{microtype}

%----------------------------- Geometría -------------------------------
\usepackage[letterpaper, margin=1in]{geometry}

\usepackage{setspace}
\setstretch{1.2}
\setlength{\emergencystretch}{3em}  % último recurso contra renglones desbordados

%---------------------------- Matemáticas ------------------------------
\usepackage{amsmath,amssymb,amsthm,mathtools}

%------------------------ Figuras, tablas, código ----------------------
\usepackage{graphicx}
\graphicspath{{figuras/}}
\usepackage{float}
\usepackage{booktabs}
\usepackage{array}
\usepackage[font=small,labelfont=bf]{caption}
\usepackage{subcaption}             % varias gráficas lado a lado en una figura
\usepackage{enumitem}
\usepackage{xcolor}
% Gráficas hechas en LaTeX: los datos van en el .tex y el modelo se evalúa al
% compilar, así que no hay que regenerar imágenes al cambiar un parámetro.
\usepackage{pgfplots}
\pgfplotsset{compat=1.18}

% Dentro de flotantes y de sus leyendas se vuelve al interlineado sencillo.
% setspace ya lo hace para los flotantes normales, pero NO para los colocados
% con [H]; este parche cubre ambos casos.  Va después de \usepackage{float}.
\makeatletter
\g@addto@macro\@floatboxreset{\def\baselinestretch{1}\normalsize}
\makeatother

%------------------- Listados de código --------------------------------
% El ejercicio 6 pide el código en un archivo aparte; si quieres mostrar
% fragmentos aquí: \lstinputlisting[firstline=..,lastline=..]{archivo.nlogo}
\usepackage{listings}
\lstdefinestyle{codigo}{
  basicstyle=\setstretch{1}\ttfamily\footnotesize,   % \setstretch va primero
  keywordstyle=\bfseries,
  commentstyle=\itshape\color{gray!70!black},
  stringstyle=\color{teal!70!black},
  numbers=left, numberstyle=\tiny, numbersep=6pt,
  frame=tb, breaklines=true, showstringspaces=false,
  captionpos=b, xleftmargin=1.5em
}
\lstset{style=codigo}
\lstset{extendedchars=true, inputencoding=utf8,
  literate=%
  {á}{{\'a}}1 {é}{{\'e}}1 {í}{{\'i}}1 {ó}{{\'o}}1 {ú}{{\'u}}1
  {Á}{{\'A}}1 {É}{{\'E}}1 {Í}{{\'I}}1 {Ó}{{\'O}}1 {Ú}{{\'U}}1
  {ñ}{{\~n}}1 {Ñ}{{\~N}}1 {ü}{{\"u}}1 {Ü}{{\"U}}1
  {¿}{{?`}}1 {¡}{{!`}}1
}
\renewcommand{\lstlistingname}{Listado}

%------------------------- Notas de trabajo ----------------------------
% \pendiente{...} marca en rojo lo que falta escribir.  Se usa vacío —\pendiente{},
% que imprime «[pendiente]»— o con una nota de QUÉ va en ese hueco, nunca de cómo
% escribirlo.  Al terminar, descomenta la última línea y todas las notas
% desaparecen sin tener que borrarlas una por una.
\newcommand{\pendiente}[1]{%
  \textcolor{red!70!black}{[\textbf{pendiente}%
    % misma prueba de «argumento vacío» que el entorno ejercicio
    \if\relax\detokenize{#1}\relax\else\textbf{:}\ #1\fi]}}
%\renewcommand{\pendiente}[1]{}

%--------------------------- Bibliografía ------------------------------
% APA 7ª edición: el enunciado ya cita en autor-año, «(Tetschke, 2018)».
% \DeclareLanguageMapping NO es opcional: sin él la lista sale con «Retrieved
% from» y demás en inglés aunque el documento esté en español.
% Se cita con \parencite{x} -> (Tetschke et al., 2018) y \textcite{x}.
\usepackage[style=apa, backend=biber]{biblatex}
\DeclareLanguageMapping{spanish}{spanish-apa}
\addbibresource{bibliografia.bib}

%--------------------------- Hipervínculos -----------------------------
% Ligas activas pero en negro, que es lo adecuado para imprimir.  hyperref se
% carga casi al final; cleveref SIEMPRE después de hyperref.
\usepackage[hidelinks]{hyperref}
\hypersetup{
  pdftitle={Biología Matemática I — Miniproyecto 1},
  pdfauthor={José Rubén Alfaro Gónzalez},
}
% Ligas en color para la etapa de escritura:
%\hypersetup{colorlinks=true, linkcolor=red!70!black,
%            citecolor=green!45!black, urlcolor=cyan!70!black}

% A cleveref hay que pasarle el idioma A MANO: no lo lee de babel, y sin esto
% escribe «Chapter 4» en pleno texto en español.  El último idioma gana.
\usepackage[english,spanish,nameinlink,noabbrev]{cleveref}
% En español cleveref usa «apartado» y «Cuadro»; aquí se prefiere el uso mexicano
\crefname{section}{secci\'on}{secciones}
\Crefname{section}{Secci\'on}{Secciones}
\crefname{table}{tabla}{tablas}
\Crefname{table}{Tabla}{Tablas}
\crefname{figure}{figura}{figuras}
\Crefname{figure}{Figura}{Figuras}

%--------------------- Ejercicios ----------------------------------------
% Va después de cleveref porque registra sus propios nombres.
%   \begin{ejercicio}             -> «Ejercicio 3»
%   \begin{ejercicio}[Tipo 1]     -> «Ejercicio 3 (Tipo 1)»
%   \begin{solucion} … \end{solucion}  cierra con un cuadrito.
% Con \label{ej:...} dentro del entorno, \Cref{ej:...} escribe «Ejercicio 3».
% Como sólo se resuelven el 1, el 2 y el 6, cada archivo fija su número con
% \setcounter antes de abrir el entorno: así conserva el del enunciado.
\newcounter{ejercicio}
\newenvironment{ejercicio}[1][]{%
  \refstepcounter{ejercicio}%
  \bigskip\par\noindent
  {\bfseries\large Ejercicio \theejercicio
   % prueba estándar de «argumento opcional vacío»: \ifx falla aquí porque
   % \newcommand define la macro como \long y \@empty no lo es.
   \if\relax\detokenize{#1}\relax\else\ (#1)\fi}%
  \par\nobreak\smallskip\noindent\ignorespaces}
  {\par\medskip}
\newenvironment{solucion}
  {\par\smallskip\noindent\textit{Soluci\'on.}\ \ignorespaces}
  {\hfill$\square$\par\medskip}
\crefname{ejercicio}{ejercicio}{ejercicios}
\Crefname{ejercicio}{Ejercicio}{Ejercicios}
